\begin{abstract}
Extremely large-scale multiple-input multiple-output (XL-MIMO) is a pivotal technology for 6G communications; however, it introduces severe near-field spatial non-stationarity. Consequently, individual users are effectively covered by only a localized subset of antennas, known as the visibility region (VR). Under these conditions, traditional full-array processing becomes computationally prohibitive and susceptible to noise accumulation, whereas fully-connected graph neural network (GNN) detectors suffer from severe scalability bottlenecks. To address these challenges, this paper proposes Graph-VR-Det, a graph-based near-field detection scheme. The proposed method utilizes pilot-energy distributions for dynamic VR discovery to construct sparse user-antenna bipartite graphs. By restricting the evidence processing of a sparse message passing neural network (Sparse-GNN) exclusively to valid edges, computational complexity is significantly reduced. Furthermore, a frequency-aware extension is introduced to compensate for wideband beam splitting effects at terahertz (THz) frequencies. Simulations demonstrate that the proposed Sparse-GNN achieves up to a 3.3-fold inference speedup over fully-connected networks. At an optimal sparsity level, the proposed method exhibits a graph regularization effect that filters out noise, surpassing the bit error rate performance of dense networks. The architecture also demonstrates zero-shot scalability to 1024-antenna arrays without retraining. However, in wideband scenarios, the proposed method exhibits an error floor at high signal-to-noise ratios compared to ideal genie-aided linear detectors, highlighting a trade-off between computational efficiency and high-SNR precision.
\end{abstract}